\subsection{A second homogeneous candidate: independent linear and cyclic orders}

\begin{proposition}[Refuted: the linear--cyclic order candidate]
\label{A:cyclic-linear}
Let $C$ be a skeleton of the category of nonempty finite sets with
an independent linear order and cyclic order, whose arrows are
injections preserving both orders. Cyclic orders on one or two
points are taken to be unique. This small category is neither total
nor exact. More precisely, $BC$ is simply connected and
$H_2(BC;\mathbb Z)\cong\mathbb Z$, and there is a five-arrow parallel
family with no affine equalizer over any nonzero commutative ring.
No assertion that $C$ is quasi-filtered is made.
\end{proposition}
\begin{proof}
Take $p$ to have two points and $q$ to have linear order
$1<2<3<4$ and cyclic order $(1,3,2,4)$. Every increasing pair
$ij$ defines an arrow $a_{ij}:p\to q$. For any target $r$, choose
a cut in its cyclic order, giving a linear reading of that circle.
For an arrow $u:p\to r$, define $\chi(u)\in\{0,1\}$ to be one
if the larger of its two points in the given linear order of $r$
occurs first in the circular reading, and zero otherwise. Extend
$\chi$ linearly to $R[C(p,r)]\to R$ for any commutative ring $R$.

For every $g:q\to r$, the circular order induced on its image is
$(g1,g3,g2,g4)$. Examining the four possible positions of the
restricted cut gives
\[
\begin{split}
 \chi(ga_{12})+\chi(ga_{23})-\chi(ga_{13})&=1,\\
 \chi(ga_{12})+\chi(ga_{24})-\chi(ga_{14})&=0.
\end{split}
\]
These identities can also be read from the negative orientation
of the triple $123$ and the positive orientation of $124$.
For $t\in R[C(q,r)]$ annihilating the differences of the five
arrows
\[
 a_{12},a_{23},a_{13},a_{24},a_{14},
\]
write $c$ for the common value $\chi(ta_{ij})$. The first identity
gives $c=\aug(t)$ and the second gives $c=0$. Thus every simultaneous
annihilator has augmentation zero, at every target. In particular
there is no affine equalizer when $1_R\ne0$. The column exact
sequences of Proposition~\ref{B:lead3}, with $R$ in place of
$\mathbb Z$, give $T_1(Q_a)\cong R$ directly, without assuming
condition I. In the integral case the smaller four-arrow family
$a_{12},a_{23},a_{34},a_{14}$ already has an obstruction:
\[
 \chi(ga_{12})+\chi(ga_{23})+\chi(ga_{34})-\chi(ga_{14})=1
\]
forces $2c=1$ for an affine equalizer.

To check totality independently, let $X_n$ be the set of cyclic
orders on the linearly ordered set $\{0,\ldots,n\}$, with face maps
given by deletion followed by increasing relabeling. This is a
semisimplicial set. Its category of simplices is exactly $C$:
an increasing injection is an arrow precisely when the cyclic
order on its source is the restriction of that on its target.
The usual simplex-category comparison gives
$BC\simeq |X|$. One may obtain this comparison by realizing the
diagram of standard simplices over the simplex category. Its
objectwise collapse to points is a homotopy-colimit equivalence,
and its augmentation to $|X|$ is the simplex-category subdivision
comparison. Equivalently, barycentric subdivision with the face
occurrences retained gives the nerve of that category. This is
the same semisimplicial comparison used in the finite sphere
test, and does not identify the unrelated full cyclic category
with $C$.

There is one zero-simplex, one one-simplex, two two-simplices,
and six three-simplices. The cellular differentials in these
degrees, with the alternating deletion convention, are
\[
 d_1=0,\qquad
 d_2=(1,1):\mathbb Z^2\longrightarrow\mathbb Z,\qquad
 d_3=0:\mathbb Z^6\longrightarrow\mathbb Z^2.
\]
For $d_2$, all three faces of either cyclic triangle are the
unique edge, with coefficient sum $1-1+1=1$. For $d_3$, the six
cyclic orders fixing the initial point have alternating
triangle-face sums zero in each of the two orientation classes.
The program \texttt{cyclic\_linear\_cuts.py} prints and checks
all six columns, as well as all four cut identities above.
It follows that $H_2(|X|;\mathbb Z)=\ker d_2\cong\mathbb Z$.
Its one-skeleton is a circle, and either two-simplex attaches
with degree one on that circle. Hence $\pi_1(|X|)=0$.
Already the constant singleton $C$-set has element category
$C$, whose nonzero second homology excludes totality.
\end{proof}
